%=====================================================================
%  3. Method
%=====================================================================
\section{The \sysname{} Framework}
\label{sec:method}

\subsection{Overview}
\label{ssec:overview}

\sysname{} formulates personalized resource generation as a coordinated
multi-agent decision process rather than a single inference call. Given a
learning topic $t$ and a learner profile $p$, the system returns a bundle
$\mathcal{R} = \{r_1,\ldots,r_M\}$ of resources across $M$ modalities together
with a verification report and an adaptation decision. Formally,

\begin{equation}
\mathcal{R}, \hat{h}, a \;=\; \mathcal{F}(t, p, \mathcal{K}),
\end{equation}

where $\mathcal{K}$ is the domain knowledge base, $\hat{h}$ is the estimated
hallucination rate of the bundle, and $a$ is the adaptation action for the next
iteration.

Figure~\ref{fig:arch} depicts the architecture. A supervisor decomposes the
request and routes it; a diagnosis agent converts the profile into an explicit
difficulty target and gap list; six specialist agents generate resources from a
shared state; a debate agent adversarially verifies the bundle; a review agent
scores it; and---only if verification passes---an adaptation agent decides the
next difficulty transition. A critical structural property is that
\emph{verification precedes delivery}: the loop back to generation is triggered
by severity-graded evidence, not by a fixed iteration count.

\begin{figure*}[!t]
\centering
%=====================================================================
%  System architecture — TikZ (included from main.tex preamble styles)
%=====================================================================
\begin{tikzpicture}[
  scale=1.1,
  ctrl/.style={draw, rounded corners=3pt, minimum width=2.8cm, minimum height=0.66cm,
               align=center, font=\scriptsize, fill=blue!10, draw=blue!50},
  gen/.style={draw, rounded corners=2pt, minimum width=1.42cm, minimum height=0.6cm,
              align=center, font=\tiny, fill=green!10, draw=green!45},
  ver/.style={draw, rounded corners=3pt, minimum width=2.6cm, minimum height=0.66cm,
              align=center, font=\scriptsize, fill=red!10, draw=red!50},
  out/.style={draw, rounded corners=3pt, minimum width=2.5cm, minimum height=0.66cm,
              align=center, font=\scriptsize, fill=gray!10, draw=gray!60},
  kb/.style={draw, rounded corners=3pt, minimum width=1.5cm, minimum height=1.9cm,
             align=center, font=\scriptsize, fill=orange!10, draw=orange!55},
  ar/.style={-Stealth, thin, draw=black!75},
  dar/.style={-Stealth, thin, dashed, draw=black!55},
]

\node[ctrl] (prof)  at (5.0,  0.0) {Learner profile $p$};
\node[ctrl] (sup)   at (5.0, -1.2) {Supervisor};
\node[ctrl] (diag)  at (5.0, -2.4) {Diagnosis agent};

\node[gen] (g1) at (0.8, -3.9) {MindMap};
\node[gen] (g2) at (2.5, -3.9) {Document};
\node[gen] (g3) at (4.2, -3.9) {Reading};
\node[gen] (g4) at (5.9, -3.9) {Code};
\node[gen] (g5) at (7.6, -3.9) {Quiz};
\node[gen] (g6) at (9.3, -3.9) {Video};

\node[ver] (deb)  at (5.0, -5.3) {Debate \& cross-validation};
\node[ver] (rev)  at (5.0, -6.5) {Quality review};
\node[out] (adp)  at (2.9, -7.9) {Adapt agent};
\node[out] (dlv)  at (7.4, -7.9) {Delivered bundle};

\node[kb]  (kb)   at (11.7, -4.7) {Domain\\knowledge\\base};

\draw[ar]  (prof) -- (sup);
\draw[ar]  (sup)  -- (diag);

\foreach \g in {g1,g2,g3,g4,g5,g6}{
  \draw[ar] (diag.south) -- ++(0,-0.45) -| (\g.north);
}

\foreach \g in {g1,g2,g3,g4,g5,g6}{
  \draw[ar] (\g.south) |- (deb.north);
}

\draw[ar]  (deb) -- (rev);
\draw[ar]  (rev.south) -- ++(0,-0.55) -| (adp.north) node[pos=0.15,above,font=\tiny]{pass};
\draw[ar]  (rev.south) -- ++(0,-0.55) -| (dlv.north);

\draw[dar] (rev.east) -- ++(0.9,0) -- ++(0,2.6) -- (g6.east)
            node[pos=0.35,above,font=\tiny]{retry (critical only)};

\draw[dar] (adp.west) -- ++(-1.9,0) -- ++(0,7.1) -- (prof.west)
            node[pos=0.6,above,font=\tiny]{profile update};

\draw[<->,thin,draw=black!70] (kb.west) -| (deb.east)
            node[pos=0.25,above,font=\tiny]{evidence};
\draw[dotted,thin,draw=black!60] (kb.north) -- (g6.east);

\end{tikzpicture}

\caption{Architecture of \sysname{}. Solid arrows denote the forward
generation path; dashed arrows denote verification-driven retry and the
feedback path that updates the learner profile across sessions.}
\label{fig:arch}
\end{figure*}

\subsection{Dialogue-Driven Learner Profiling}
\label{ssec:profile}

Rather than requiring learners to complete a survey, the profile agent
extracts features incrementally from natural conversation. The profile is a
structured record $p = \langle e, j, s, w, g, c, r \rangle$ whose fields are
listed below, with $s$ denoting per-subject theory test scores in $[0,100]$:

\begin{itemize}
\item \textbf{$e$ educational background} (e.g., vocational diploma, BSc, MSc,
PhD) and \textbf{$j$ job target} (in-school, career changer, in-service
upskilling, corporate trainee);
\item \textbf{$s$ theory test results} per subject, used as the quantitative
basis for difficulty calibration;
\item \textbf{$c$ cognitive style} (visual / auditory / kinesthetic /
read-write);
\item \textbf{$w$ identified weaknesses}, each with a subject and a severity;
\item \textbf{$g$ learning goal}, \textbf{$r$ preferred pace}, interest
preferences, and motivation level.
\end{itemize}

Because extraction is incremental, the profile is refreshed whenever new
evidence appears (a quiz result, an explicit statement, a resource
completion event), which is what makes the closed loop in
Figure~\ref{fig:arch} meaningful across sessions rather than within a single
request.

\subsection{Diagnosis and Difficulty Calibration}
\label{ssec:diagnosis}

The diagnosis agent maps $p$ onto an explicit difficulty target. Let
$\mathcal{S}(t)$ be the set of subjects a topic $t$ draws on. The base
competence score is

\begin{equation}
b(p,t) = \left\lfloor \frac{1}{|\mathcal{S}(t)|}
\sum_{u \in \mathcal{S}(t)} s_u + \frac{1}{2} \right\rfloor,
\end{equation}

where a missing subject score is imputed by the median of the learner's
observed scores, rounded by the same half-up rule. We adopt round-half-up
rather than a ceiling so that the base score is unbiased with respect to the
underlying per-subject scores; the implementation
(\texttt{orchestrator/workflow.py}, \texttt{tests/evaluate\_quality.py}) and the
dataset specification (\texttt{test\_cases.json}) use this identical rule, so
the expected label is reproducible across the system, the harness, and a
reviewer's own recomputation. The calibrated target score is then

\begin{equation}
d(p,t) = \mathrm{clip}_{[0,100]}\bigl(b(p,t) + \delta_g + \delta_r + \delta_j\bigr),
\end{equation}

with $\delta_g, \delta_r, \delta_j \in \mathbb{Z}$ the goal, pace, and
career-stage modifiers. The difficulty label is
$\ell = \textsc{basic}$ if $d < 60$, $\textsc{intermediate}$ if $60 \le d < 85$,
and $\textsc{advanced}$ otherwise. Crucially, this rule is \emph{public and
deterministic}: given a profile and a topic, any reviewer can recompute the
expected label, which is what makes the adaptation accuracy reported in
Section~\ref{sec:results} auditable instead of self-reported.

The computed quantities---prerequisite subjects, aggregate score $b$, modifier
terms, and resulting level $\ell$---are injected into the supervisor's
instruction together with the boundaries, so that the decision is a checked
computation rather than an impression. The supervisor retains an escape hatch:
it may deviate from the computed level only by stating concrete evidence, and
the deviation is logged. Our first measurement showed why this matters: without
the injected rule, the same supervisor produced a systematic one-level-downward
bias (Section~\ref{ssec:calib_ablation}).

\subsection{Specialized Generation Agents}

Six specialists operate on a shared state object, so that later agents observe
what earlier ones produced. This shared state is the mechanism behind
cross-modal consistency: the quiz agent writes questions about the same
knowledge points the document agent covered, and the video agent scripts
against the same outline the mind map produced.

\begin{table}[!t]
\centering
\caption{Agents participating in the collaboration workflow.}
\label{tab:agents}
\begin{tabular}{@{}ll@{}}
\toprule
\textbf{Agent} & \textbf{Output} \\
\midrule
Supervisor & Task decomposition, agent routing, difficulty plan \\
Diagnosis & Gap localization, difficulty target, priority focus \\
MindMap & Hierarchical knowledge structure reused downstream \\
Document & Personalized explanatory prose at target difficulty \\
Reading & Academic frontier extension material \\
Code & Runnable, commented code case \\
Quiz & Graded questions (5 item types) with rubric \\
Video & Lecture video script with scene breakdown \\
Debate & Hallucination findings, cross-agent contradiction rulings \\
Review & Four-dimension quality scores and issue list \\
Adapt & Difficulty transition decision for next iteration \\
\bottomrule
\end{tabular}
\end{table}

\subsection{Debate and Cross-Validation}
\label{ssec:debate}

The debate agent is the component that most directly addresses the reliability
problem. It receives the union of generated resources and performs two
distinct checks. \emph{Internal cross-validation} looks for contradictions
between agents---for example, a code sample using an API signature that the
document describes differently---and issues an explicit ruling for each.
\emph{External fact-checking} compares each extracted key claim against
retrieved knowledge-base evidence and returns, per claim, a verdict with a
pointer to supporting text.

Findings are graded by severity. Let $\mathcal{C}$ be the set of checked
claims and $\mathcal{H} \subseteq \mathcal{C}$ the subset flagged as
unsupported. The bundle-level hallucination rate is

\begin{equation}
\hat{h} = \frac{|\mathcal{H}|}{|\mathcal{C}|}.
\end{equation}

Algorithm~\ref{alg:verify} gives the control flow. A retry is triggered only
when a finding is graded \textsc{critical} \emph{and} the verifier recommends
one, and at most one retry is permitted; this bounded policy prevents the
unbounded self-critique loops that plague iterative refinement methods.

\begin{algorithm}[!t]
\caption{Verification-gated generation (per request)}
\label{alg:verify}
\begin{algorithmic}[1]
\STATE $p \leftarrow$ profile agent extraction or lookup
\STATE $\tau \leftarrow$ diagnosis$(p, t)$ \COMMENT{difficulty target and gaps}
\STATE $\mathcal{R} \leftarrow \emptyset$
\FOR{$a \in \{\text{MindMap}, \text{Doc}, \text{Reading}, \text{Code}, \text{Quiz}, \text{Video}\}$}
  \STATE $\mathcal{R} \leftarrow \mathcal{R} \cup \{ a(\tau, \mathcal{K}, \mathcal{R}) \}$ \COMMENT{shared state}
\ENDFOR
\FOR{$i \in \{1,2\}$} \COMMENT{at most one retry}
  \STATE $(\mathcal{H}, \hat{h}) \leftarrow$ debate$(\mathcal{R}, \mathcal{K})$
  \STATE $(q, \mathcal{I}) \leftarrow$ review$(\mathcal{R}, \tau)$
  \IF{$\neg \text{critical}(\mathcal{H}) \;\wedge\; \neg \text{critical}(\mathcal{I})$}
    \STATE \textbf{break}
  \ELSIF{$i = 1$}
    \STATE $\mathcal{R} \leftarrow$ regenerate$(\tau, \mathcal{K}, \mathcal{H}, \mathcal{I})$
  \ENDIF
\ENDFOR
\STATE $a \leftarrow$ adapt$(p, \tau, \mathcal{R}, q)$ \COMMENT{next-iteration decision}
\STATE \textbf{return} $(\mathcal{R}, \hat{h}, q, a)$
\end{algorithmic}
\end{algorithm}

\subsection{Quality Review and Adaptive Decision}

The review agent scores the bundle on four dimensions---accuracy, difficulty
fit, completeness, and cross-modal consistency---each on a 10-point scale, and
emits a severity-graded issue list. The adaptation agent then applies an
explicit, publishable rule (Table~\ref{tab:adapt}) that maps mastery and
observed accuracy onto one of four actions. Stating this rule explicitly is
deliberate: it converts ``the model adapts'' from an unfalsifiable claim into a
testable one.

\begin{table}[!t]
\centering
\caption{Adaptation decision rules. $\mu$ is mastery on the knowledge point,
$\alpha$ quiz accuracy, $q_{\text{fit}}$ the reviewer's difficulty-fit score.}
\label{tab:adapt}
\begin{tabular}{@{}ll@{}}
\toprule
\textbf{Condition} & \textbf{Action} \\
\midrule
$\mu < 60$ \textbf{or} $q_{\text{fit}} < 6$ & Dimension reduction \\
$\mu \ge 85$ \textbf{and} $\alpha \ge 80\%$ & Advanced challenge \\
critical issue pending & Content correction \\
otherwise & Maintain \\
\bottomrule
\end{tabular}
\end{table}

\subsection{Retrieval Grounding}

Domain knowledge is stored as chunked documents embedded in a vector store.
Every generation agent receives retrieved context, and the debate agent reuses
the same index when checking claims, so generation and verification are
grounded in \emph{the same} evidence rather than in two inconsistent views of
the domain.

\subsection{Implementation}

The orchestration layer is implemented as an explicit state graph; the service
layer exposes it over HTTP with server-sent events for incremental display; the
application layer (learner management, quiz administration, persistence) is a
conventional web service. All agents in the reported experiments call the same
commercial LLM endpoint with per-agent temperature settings, so no component
has access to a differently-capable model.
